\chapter{Arithmetic Positselski: Primitive Zeta-Character Theorem}
\label{v4-arith:prim-zeta-char}

The finite primary towers have an ordinary trace \(\zeta(s)\) in the
half-plane \(\operatorname{Re}s>1\). The Gaussian test function has a
Mellin integral \(\Gamma_{\mathbb R}(s)\). Their product gives the
completed zeta function, with the two operations kept distinct. The
normalized thermal state on the finite tower is a further functional,
with unit value one. We compute these objects and specify the maps
that would transfer their scalar identity to a native Koszul character.

\section{Statement}
\label{v4-arith:prim-zeta:statement}

\begin{definition}[mixed trace and Mellin character]
\label{v4-arith:prim-zeta:mixed-character}
Let \(P\) be the finite-support occupation space with basis \(v_m\),
energy \(Hv_m=(\log m)v_m\), and Hilbert completion \(\ell^2(\mathbb N^\times)\).
For the specified Gaussian \(g(x)=e^{-\pi x^2}\), define
\[
 \chi_{g,P}(s)=\left(\int_{\mathbb R}g(x)|x|^{s-1}dx\right)
                    \operatorname{Tr}_{\ell^2}(e^{-sH}),
                    \qquad\operatorname{Re}s>1.
\]
The symbol \(\mathrm{char}\) for the completed primary scalar below
means this mixed operation. Its finite factor is an ordinary graded
trace; its archimedean factor is a Mellin functional on a test function.
\end{definition}

\begin{hypothesis}[native character realization]
\label{v4-arith:prim-zeta:native-character-hypothesis}
To identify the mixed scalar with characters of native bar and cobar
objects, assume those objects and their graded trace functionals are
constructed. Assume specified graded linear comparison maps identify
their trace-carrying sectors with the represented occupation carrier,
using the trace parameter \(s\) and \(1-s\), respectively, and identify their archimedean
functionals with the Gaussian Mellin functional at \(s\) and at
\(1-s\). Each map must preserve the relevant trace and energy grading
on a common convergence domain. A claim about the full native object
also requires that all complementary sectors have zero contribution
for that trace. The comparison of multiplications and higher operations
is an additional requirement, separate from this trace hypothesis.
\end{hypothesis}

\begin{theorem}[primitive zeta-character theorem]
\label{v4-arith:prim-zeta:thm:main}
\ClaimStatusProvedHere. Let
\(\mathcal{V}^{\mathrm{prim}}_{\mathrm{Heis}}=\bigotimes'_{p}\mathcal{H}^{(p)}\otimes\mathcal{H}^{(\infty)}\)
be the adelic arithmetic Heisenberg Fock and let
\(\mathcal{V}^{\mathrm{prim}}_{\mathrm{Heis},\mathrm{prim}}\) be its
\emph{primary subspace} as defined in
\Cref{v4-arith:prim-zeta:def:primary-subspace}. Endow both with the
HY-zeta grading of
\Cref{v4-arith:conv:two-grading}(ii). Then the mixed trace and Mellin character of
the primary datum is the completed zeta function:
\begin{equation}
\label{v4-arith:prim-zeta:eq:main}
\mathrm{char}\bigl(\mathcal{V}^{\mathrm{prim}}_{\mathrm{Heis},\mathrm{prim}}\bigr)(s)
\;=\;\Gamma_{\mathbb{R}}(s)\,\zeta(s)
\;=\;\Lambda_{\zeta}(s),
\qquad
\Gamma_{\mathbb{R}}(s)\;:=\;\pi^{-s/2}\Gamma(s/2).
\end{equation}
Under \Cref{v4-arith:prim-zeta:native-character-hypothesis}, the
character of the Koszul dual, under the arithmetic bar
coalgebra \(C=B^{\mathrm{ch,Ar}}(\mathcal{V}^{\mathrm{prim}}_{\mathrm{Heis},\mathrm{prim}})\)
and cobar algebra \(A=C^{\vee}\), satisfies:
\begin{equation}
\label{v4-arith:prim-zeta:eq:koszul-dual}
\mathrm{char}_{C}(s)\;=\;\Lambda_{\zeta}(s),
\qquad
\mathrm{char}_{A^{!}}(s)\;=\;\Lambda_{\zeta}(1-s),
\qquad
\mathrm{char}_{C}(s)\;=\;\mathrm{char}_{A^{!}}(1-s).
\end{equation}
In particular, the functional equation
\(\Lambda_{\zeta}(s)=\Lambda_{\zeta}(1-s)\) (Tate 1950)
gives \(\mathrm{char}_{C}(s)=\mathrm{char}_{A^{!}}(s)\) as an identity
of meromorphic functions.
The mixed scalar identity is unconditional on its convergence domain
and extends meromorphically, with poles at zero and one. The native
character identities require
\Cref{v4-arith:prim-zeta:native-character-hypothesis}; neither a scalar
functional equation nor local thermal normalization constructs those
comparison maps.
\end{theorem}

\begin{remark}[range of the scalar construction]
\label{v4-arith:prim-zeta:rem:domain}
The occupation trace and Gaussian Mellin integral are explicit. To
interpret their product as a native arithmetic bar or cobar character,
the trace maps of \Cref{v4-arith:prim-zeta:native-character-hypothesis}
remain necessary. Equality of scalar functions does not identify
native graded objects, their coefficient modules or their products.
\end{remark}

\begin{remark}[why the primary subspace is the right object]
\label{v4-arith:prim-zeta:rem:why-primary}
The full arithmetic Heisenberg Fock
\(\mathcal{V}^{\mathrm{prim}}_{\mathrm{Heis}}\) carries character
\(\prod_{p}\prod_{n}(1-p^{-ns})^{-1}=\prod_{n\geq 1}\zeta(ns)\) under
HY-zeta grading, not \(\Lambda_{\zeta}(s)\). The factor
\(\prod_{n\geq 2}\zeta(ns)\) beyond \(\zeta(s)\) encodes the
multi-oscillator Fock sector and produces higher-weight Dirichlet
series with their own trivial zeros; it is the Fock-symmetric
completion of \(\zeta(s)\), not \(\zeta(s)\) itself. The archimedean
Fock \(\mathcal{H}^{(\infty)}\) contributes a product
\(\prod_{n\geq 1}(1-e^{-n/T})^{-1}\) under temperature \(T\) conjugate to
conformal weight; its Mellin transform at criticality is
\(\Gamma(s)\zeta(s)\), not \(\Gamma_{\mathbb{R}}(s)\zeta(s)\).

The primary subspace is defined to extract (i) the single-oscillator
sector at each prime (which gives \(\zeta(s)\) not
\(\prod_{n\geq 1}\zeta(ns)\)), and (ii) the Gaussian-primary archimedean
sector (which gives \(\Gamma_{\mathbb{R}}(s)\) not \(\Gamma(s)\)). The
tensor gives \(\Lambda_{\zeta}(s)\).
\end{remark}

\section{The finite trace and the Gaussian Mellin integral}
\label{v4-arith:prim-zeta:table}

The finite trace and the archimedean integral use different carriers.
Their product and its functional equation are scalar identities; native
character transfer uses the additional comparison hypothesis.

\subsection{The Finite and Archimedean Primary Object
\texorpdfstring{\(\mathcal{V}^{\mathrm{prim}}_{\mathrm{Heis},\mathrm{prim}}\)}{V\^{}prim\_{Heis,prim}}}
\label{v4-arith:prim-zeta:att-1-primary}

\paragraph{Definition.} The notation
\(\mathcal{V}^{\mathrm{prim}}_{\mathrm{Heis},\mathrm{prim}}\) has been
used in
\Cref{v4-arith:vb:def:T_n},
\Cref{v4-arith:rh:conj:primary-Virasoro}, and
\Cref{v4-arith:rh:prop:Li-is-Virasoro} as ``the primary subspace of
the arithmetic Heisenberg Fock''.  Here it is the adelic object whose
finite local factors have character \((1-p^{-s})^{-1}\) and whose
archimedean factor is the Gaussian Tate line.

\paragraph{Finite places.} At a finite prime \(p\), the primary factor
is the polynomial tower generated by the single line
\(\mathbb{C}\cdot a^{(p)}_{-1}|0\rangle_{p}\).  Its basis is
\(\{e_{n,p}\}_{n\geq0}\), where \(e_{0,p}=|0\rangle_{p}\) and
\(L_{0}^{\mathrm{Ar}}e_{n,p}=n\log p\,e_{n,p}\).  The finite vacuum is
the local unit of the restricted tensor product. It is not a
trivial-zero state.

\paragraph{Archimedean place.} At \(\infty\), the primary factor is
the Gaussian line \(\mathbb{C}\cdot e^{-\pi x^{2}}\).  Its Mellin
pairing over \(\mathbb{R}^{\times}\) is
\(\Gamma_{\mathbb R}(s)\).  The Hermite tower above the Gaussian is
not part of the primary line; that tower is the archimedean
trivial-zero sector.

\paragraph{Character.}
\(\mathcal{V}^{\mathrm{prim}}_{\mathrm{Heis},\mathrm{prim}}\) is
the restricted tensor product of these factors.  Its finite character
is \(\prod_{p}(1-p^{-s})^{-1}=\zeta(s)\); its archimedean character is
\(\Gamma_{\mathbb R}(s)\).

\subsection{Mellin--Theta Character Computation}
\label{v4-arith:prim-zeta:att-2-mellin}

\paragraph{Computation.} Compute the per-place partition function
separately and take a restricted tensor product. At each finite
\(p\), the primary contribution is \((1-p^{-s})^{-1}\) (single
oscillator). Product: \(\zeta(s)\). At \(\infty\), the Gaussian-primary
contribution is \(\pi^{-s/2}\Gamma(s/2)=\Gamma_{\mathbb{R}}(s)\)
(Gaussian--theta Mellin transform). Tensor:
\(\Gamma_{\mathbb{R}}(s)\zeta(s)=\Lambda_{\zeta}(s)\).

\paragraph{Archimedean factor.} The archimedean factor is exactly
\(\Gamma_{\mathbb{R}}(s)\), not \(\Gamma(s)\). Tate 1950 \S2.5:
the archimedean local \(L\)-factor for the trivial Dirichlet character
on \(\mathbb{Q}_{\infty}=\mathbb{R}\) is
\(L_{\infty}(s)=\pi^{-s/2}\Gamma(s/2)=\Gamma_{\mathbb{R}}(s)\);
the archimedean theta vector is \(e^{-\pi x^{2}}\), which is
Gaussian, and its Mellin transform against \(|x|^{s-1}\) is exactly
\(\Gamma_{\mathbb{R}}(s)/2\). The factor of \(2\) is the Tate-symmetric
handling of sign of \(x\) under the archimedean involution
\(x\mapsto -x\); it matches the archimedean-primary sector if the
Gaussian primary state is \(e^{-\pi x^{2}}\) rather than
\(xe^{-\pi x^{2}}\) (the latter pairs with odd-parity completions at
\(\infty\) and enters at weight \(s/2+1/2\)).

\paragraph{Result.} Primary archimedean Gaussian
\(\to\Gamma_{\mathbb{R}}(s)\) exactly. Finite-place primary
single-oscillator \(\to\prod_{p}(1-p^{-s})^{-1}=\zeta(s)\). Adelic
tensor: \(\Lambda_{\zeta}(s)\), by Tate 1950 Mellin--theta
(\Cref{v4-arith:cl:P2-resolution}) and Julia 1990 primon gas.

\subsection{Koszul-Dual Character}
\label{v4-arith:prim-zeta:att-3-koszul-dual}

The substitution \(s\mapsto1-s\) defines the dual scalar
\(\chi_{g,P}(1-s)=\Gamma_{\mathbb R}(1-s)\zeta(1-s)\), initially on
\(\operatorname{Re}s<0\). Its finite factor is a convergent ordinary
trace there. It is a Gibbs partition function only when \(1-s\) is
real and greater than one. Calling it a native Koszul-dual character
requires \Cref{v4-arith:prim-zeta:native-character-hypothesis}.

\subsection{The scalar functional equation}
\label{v4-arith:prim-zeta:att-4-match}
The completed functional equation compares the meromorphic continuations
\(\chi_{g,P}(s)=\chi_{g,P}(1-s)\). It does not compare two simultaneous
Gibbs states: their displayed absolute-convergence half-planes are
disjoint. Under the native character hypothesis, it transfers to the
stated native trace functionals.

\subsection{The operations and their domains}
\label{v4-arith:prim-zeta:table-summary}
\begin{table}[h]
\centering
\begin{tabular}{lll}
Operation & Value & Initial domain\\
\hline
Occupation trace & \(\zeta(s)\) & \(\operatorname{Re}s>1\)\\
Gaussian Mellin integral & \(\Gamma_{\mathbb R}(s)\) & \(\operatorname{Re}s>0\)\\
Mixed character & \(\Lambda_\zeta(s)\) & \(\operatorname{Re}s>1\)\\
Dual mixed character & \(\Lambda_\zeta(1-s)\) & \(\operatorname{Re}s<0\)
\end{tabular}
\caption{The scalar operations used in the primary character.}
\label{v4-arith:prim-zeta:tab:table}
\end{table}

\section{Primary Subspace
\texorpdfstring{\(\mathcal{V}^{\mathrm{prim}}_{\mathrm{Heis},\mathrm{prim}}\)}{V\^{}prim\_{Heis,prim}}}
\label{v4-arith:prim-zeta:primary-subspace}

\subsection{Local Primary Subspaces}
\label{v4-arith:prim-zeta:local-primary}

\begin{definition}[per-prime primary subspace]
\label{v4-arith:prim-zeta:def:p-primary}
\ClaimStatusDefinitional. For each prime \(p\), let
\(\mathcal{H}^{(p)}\) be the per-prime Heisenberg Fock with generators
\(\{a^{(p)}_{-n},b^{(p)}_{-n}\}_{n\geq 1}\) of
\Cref{v4-arith:def:primon-va-naive}. The \emph{per-prime primary
subspace}
\(\mathcal{H}^{(p)}_{\mathrm{prim}}\subset\mathcal{H}^{(p)}\) is
the polynomial tower on the single oscillator \(a^{(p)}_{-1}\):
\begin{equation}
\label{v4-arith:prim-zeta:eq:p-primary}
\mathcal{H}^{(p)}_{\mathrm{prim}}\;:=\;
\bigoplus_{n\geq 0}\mathbb{C}\cdot (a^{(p)}_{-1})^{n}|0\rangle_{p}
\;=\;\bigoplus_{n\geq 0}\mathbb{C}\cdot e_{n,p},
\end{equation}
where \(e_{0,p}=|0\rangle_{p}\) and
\(e_{n,p}:=(a^{(p)}_{-1})^{n}|0\rangle_{p}\) has
\(L_{0}^{\mathrm{Ar}}\)-weight \(n\log p\). The finite vacuum is the
spherical unit of the restricted tensor product. It supplies the
\(1\) in \((1-p^{-s})^{-1}\). The trivial-zero sector is not this
finite unit; it is the archimedean Hermite tower excluded in
\Cref{v4-arith:prim-zeta:def:arch-primary}. Oscillator modes
\(a^{(p)}_{-m}\) with \(m\geq2\), and products involving them, are
outside the primary factor.
\end{definition}

\begin{remark}[why only one finite oscillator remains]
\label{v4-arith:prim-zeta:rem:exclude-doubles}
The per-prime Fock character including all multi-oscillator states is
\begin{equation}
\mathrm{char}\bigl(\mathcal{H}^{(p)}\bigr)(s)\;=\;\prod_{n\geq 1}\bigl(1-p^{-ns}\bigr)^{-1},
\end{equation}
which is the primon gas local factor and carries contributions from
\(p^{-ms}\) for every oscillator mode \(m\geq1\). The local primary
keeps the \(m=1\) mode and all its occupation numbers. Its basis is
\((a^{(p)}_{-1})^{n}|0\rangle_{p}\), \(n\geq0\), and its character is
\((1-p^{-s})^{-1}\), giving \(\prod_{p}(1-p^{-s})^{-1}=\zeta(s)\). The
factor \(\prod_{m\geq2}(1-p^{-ms})^{-1}\) not in the primary factor is
the finite Fock-symmetric completion.  The trivial zeros of the
completed zeta function come from the archimedean gamma factor and
the functional equation, not from deleting the finite vacuum.
\end{remark}

\subsection{Hilbert normalization and prime creation}
\label{v4-arith:prim-zeta:hilbert-normalization}

The polynomial basis $e_{n,p}$ fixes the graded vector space. For
its analytic realization, equip this tower with the Hermitian
form $\langle e_{n,p},e_{m,p}\rangle=\delta_{nm}n!$ and put
$v_{n,p}=e_{n,p}/\sqrt{n!}$. The $v_{n,p}$ are orthonormal.
This convention concerns the displayed primary tower; it does
not construct a positive form on the full arithmetic vertex or
Hochschild object.

Multiplication by its polynomial generator sends
$v_{n,p}$ to $\sqrt{n+1}v_{n+1,p}$. The bounded prime creation
operator is instead its unit shift
\[
 S_pv_{n,p}=v_{n+1,p},\qquad
 S_pe_{n,p}=\frac{e_{n+1,p}}{\sqrt{n+1}}.
\]
Its adjoint $R_p$ kills the vacuum and lowers positive occupation
by one. Thus $\mu_{p^r}=S_p^r$ represents multiplication of the
integer index by $p^r$. The operator retains the energy increment
$r\log p$ and satisfies $R_pS_p=1$. These are the prime operators
of \Cref{v4-arith:thermal:chapter}.

\begin{proposition}[prime Gibbs normalization]
\label{v4-arith:prim-zeta:prop:BC-ergodic}
\ClaimStatusProvedHere. At a real inverse temperature $\beta>0$,
put $q=p^{-\beta}$. On the completed primary tower let
$H_pv_{n,p}=n\log p\,v_{n,p}$ with its maximal self-adjoint
domain. The unnormalized partition trace and normalized state are
\[
 Z_p(\beta)=\operatorname{Tr}(e^{-\beta H_p})=(1-q)^{-1},
 \qquad \varphi_{p,\beta}(a)=(1-q)
                 \operatorname{Tr}(a e^{-\beta H_p}).
\]
For the rank-one projection $E_{nn}$ and the tail projection
$Q_n=S_p^nR_p^n$, their expectations are
\[
 \varphi_{p,\beta}(E_{nn})=(1-q)q^n,
 \qquad \varphi_{p,\beta}(Q_n)=q^n.
\]
In particular,
$\varphi_{p,\beta}(R_pS_p)=1$ and
$\varphi_{p,\beta}(S_pR_p)=q$.
The value $q^n$ in the character is a Boltzmann weight; a state
acts on an observable, not on the basis vector $e_{n,p}$.
\end{proposition}

\begin{proof}
The operator $e^{-\beta H_p}$ is diagonal with summable
entries $q^n$. The rank-one projection selects one entry, whereas
$Q_n$ selects the geometric tail starting at $n$. This gives both
expectations and normalization at the identity. More generally,
\Cref{v4-arith:thermal:prop:thermal-values} proves
\[
 \varphi_{p,\beta}(S_p^iR_p^j)=\delta_{ij}q^i,
 \qquad
 \varphi_{p,\beta}(ab)
    =\varphi_{p,\beta}(b\alpha_{i\beta}(a)),
\]
with $\alpha_t(S_p)=p^{it}S_p$. The proof uses the actual
normal-word multiplication, so this is the KMS identity with
its sign and analytic core specified. It is not ordinary cyclicity:
$[R_p,S_p]=E_{00}$ has positive expectation $1-q$.
\end{proof}

The Hall comparison preserves this normalization on its stated
carrier. In the ordinary Jordan shuffle algebra, write
$\eta_n=1$ in dimension $n$. Its product is
$\eta_a*\eta_b=\binom{a+b}{a}\eta_{a+b}$, so the multiplicative
basis is $w_n=n!\eta_n$. Declaring $w_n$ orthonormal makes
\[
 v_{n,p}\longmapsto w_n
\]
a unitary map which sends $S_p$ to Hall multiplication by $w_1$
and $R_p$ to the backward shift. The full proof, including phases,
is \Cref{v4-arith:thermal:thm:hall-module}. The map is a map of
representations; it does not identify oscillator multiplication
on the raw $e_{n,p}$ with the bounded prime shift.

\begin{proposition}[global shift expectations]
\label{v4-arith:prim-zeta:global-shift-state}
\ClaimStatusProvedHere. On $\ell^2(\mathbb N_{\geq1})$, let
$Hv_k=\log k\,v_k$ and $\mu_nv_k=v_{nk}$. For real $\beta>1$,
\[
 \varphi_\beta(a)=\zeta(\beta)^{-1}
                    \operatorname{Tr}(a e^{-\beta H})
\]
is a normalized state. For positive integers $n,m$,
\begin{equation}
 \varphi_\beta(\mu_n\mu_m^*)=\delta_{nm}n^{-\beta},
 \qquad
 \varphi_\beta(\mu_n^*\mu_m)=\delta_{nm}.
 \label{v4-arith:prim-zeta:global-shift-expectations}
\end{equation}
\end{proposition}

\begin{proof}
If $n\ne m$, either product has zero diagonal matrix entries.
For $n=m$, the first product is the projection onto the multiples
of $n$. Its unnormalized trace is
$\sum_{r\geq1}(nr)^{-\beta}=n^{-\beta}\zeta(\beta)$.
The second product is the identity. Dividing by the partition
trace gives both formulas, including $\varphi_\beta(1)=1$.
All sums are absolutely convergent for $\beta>1$.
\end{proof}

The finite-prime tensor shift algebras have compatible normalized
states for every real $\beta>0$. Their genuine arithmetic phase
extensions do not assemble through the represented singleton
phase quotients: \Cref{v4-arith:thermal:prop:two-prime-phase}
computes the obstruction already at $\{3\}\subset\{3,2\}$.
At the same time, no global nuclear density exists at $\beta=1$
in the occupation representation, because $\sum_{k\geq1}k^{-1}$
diverges. An adelic phase or categorical comparison must therefore
be constructed separately from this character computation.

\subsection{Archimedean Primary Subspace}
\label{v4-arith:prim-zeta:archimedean-primary}

\begin{definition}[archimedean primary Heisenberg]
\label{v4-arith:prim-zeta:def:arch-primary}
\ClaimStatusDefinitional. The archimedean Heisenberg factor
\(\mathcal{H}^{(\infty)}\) is the Fock representation of the rank-\(1\)
Heisenberg vertex algebra over \(\mathbb{R}\), with creation operators
\(\{\alpha_{-n}\}_{n\geq 1}\) and vacuum \(|0\rangle_{\infty}\). The
\emph{archimedean primary subspace}
\(\mathcal{H}^{(\infty)}_{\mathrm{prim}}\subset\mathcal{H}^{(\infty)}\)
is the subspace spanned by the Gaussian primary vector
\(|\text{Gauss}\rangle=e^{-\pi x^{2}}\cdot|0\rangle_{\infty}\) under the
coordinate realization of \(\mathcal{H}^{(\infty)}\) as
\(L^{2}(\mathbb{R})\) (Segal 1981 Prop.~3.2; Frenkel--Ben-Zvi 2004
\S3.5): concretely,
\begin{equation}
\label{v4-arith:prim-zeta:eq:arch-primary}
\mathcal{H}^{(\infty)}_{\mathrm{prim}}\;:=\;\mathbb{C}\cdot|\text{Gauss}\rangle\;=\;\mathbb{C}\cdot e^{-\pi x^{2}}|0\rangle_{\infty}.
\end{equation}
Descendant states \((\alpha_{-n})^{k}|\text{Gauss}\rangle\) for
\(n\geq 1,k\geq 1\) are excluded from the archimedean primary
subspace; these correspond to the Hermite-polynomial tower above the
Gaussian and carry the trivial-zero sector of \(\xi_{\mathbb R}\).
\end{definition}

\begin{proposition}[archimedean primary character]
\label{v4-arith:prim-zeta:prop:arch-char}
\ClaimStatusProvedHere. For \(\operatorname{Re}s>0\), the Gaussian
Mellin functional on the positive half-line is
\begin{equation}
\label{v4-arith:prim-zeta:eq:arch-char}
\mathcal M_+(g)(s)
\;=\;\int_{0}^{\infty}e^{-\pi x^{2}}\cdot x^{s-1}dx
\;=\;\tfrac{1}{2}\Gamma_{\mathbb{R}}(s)
\;=\;\tfrac{1}{2}\pi^{-s/2}\Gamma(s/2).
\end{equation}
The factor of \(2\) accommodates the archimedean sign involution
\(x\mapsto -x\): extending the integral over \(\mathbb{R}\) (not just
\((0,\infty)\)) gives \(\Gamma_{\mathbb{R}}(s)\) directly, which is the
Tate normalisation.
\end{proposition}

\begin{proof}
The Gaussian primary state \(|\text{Gauss}\rangle=e^{-\pi x^{2}}\)
under Mellin pairing against \(x^{s-1}\) on \((0,\infty)\):
\begin{equation}
\int_{0}^{\infty}e^{-\pi x^{2}}\cdot x^{s-1}dx
\;=\;\tfrac{1}{2}\pi^{-s/2}\Gamma(s/2),
\end{equation}
by substitution \(u=\pi x^{2}\), \(du=2\pi x\,dx\):
\begin{equation}
\tfrac{1}{2}\int_{0}^{\infty}e^{-u}\cdot(u/\pi)^{(s-1)/2}\cdot(u/\pi)^{-1/2}du/\pi^{1/2}
\;=\;\tfrac{1}{2}\pi^{-s/2}\int_{0}^{\infty}e^{-u}u^{s/2-1}du
\;=\;\tfrac{1}{2}\pi^{-s/2}\Gamma(s/2).
\end{equation}
Extension to \(\mathbb{R}\) doubles this: the full
Mellin--theta presentation is \(\int_{-\infty}^{\infty}e^{-\pi x^{2}}|x|^{s-1}dx=\pi^{-s/2}\Gamma(s/2)=\Gamma_{\mathbb{R}}(s)\)
(Tate 1950 \S2.5; Cassels--Fr\"ohlich 1967 \S13.3).
\end{proof}

\begin{proposition}[the Gaussian line is not a gamma trace]
\label{v4-arith:prim-zeta:gaussian-line-obstruction}
No fixed linear endomorphism \(L\) of the one-dimensional Gaussian
line has \(\operatorname{Tr}(e^{-sL})=\Gamma_{\mathbb R}(s)\) on
\(\operatorname{Re}s>0\).
\end{proposition}
\begin{proof}
Write \(L=\lambda\). The proposed equality at \(s=1\) gives
\(e^{-\lambda}=\Gamma_{\mathbb R}(1)=1\). At \(s=2\) it would give
\(1=e^{-2\lambda}=\Gamma_{\mathbb R}(2)=\pi^{-1}\), a contradiction.
Thus the Gaussian Mellin functional cannot be silently replaced by a
one-dimensional heat trace.
\end{proof}

\subsection{Adelic Primary Subspace}
\label{v4-arith:prim-zeta:adelic-primary}

\begin{definition}[primary subspace
\(\mathcal{V}^{\mathrm{prim}}_{\mathrm{Heis},\mathrm{prim}}\)]
\label{v4-arith:prim-zeta:def:primary-subspace}
\ClaimStatusDefinitional. The \emph{primary subspace} of the adelic
arithmetic Heisenberg Fock is the restricted tensor product
\begin{equation}
\label{v4-arith:prim-zeta:eq:adelic-primary}
\mathcal{V}^{\mathrm{prim}}_{\mathrm{Heis},\mathrm{prim}}\;:=\;
\mathcal{H}^{(\infty)}_{\mathrm{prim}}\otimes_{\mathbb{C}}\bigotimes_{p\text{ prime}}^{\prime}\mathcal{H}^{(p)}_{\mathrm{prim}}
\end{equation}
where:
\begin{itemize}
\item At \(\infty\), \(\mathcal{H}^{(\infty)}_{\mathrm{prim}}\) is the
Gaussian primary line
(\Cref{v4-arith:prim-zeta:def:arch-primary}).
\item At each finite \(p\), \(\mathcal{H}^{(p)}_{\mathrm{prim}}\) is
the single-oscillator primary space
(\Cref{v4-arith:prim-zeta:def:p-primary}).
\item The restricted tensor product is taken with respect to the
spherical vector \(e_{0,p}:=|0\rangle_{p}\) at almost all primes in
the primary factor (Tate 1950 \S2.3; Weil 1967 Ch.~VII \S1).
Thus only finitely many finite primes carry \(e_{n,p}\) with
\(n>0\) in an elementary tensor.
\end{itemize}
The underlying \(\Lambda\)-module structure at each finite \(p\) is
free of rank one
(the Iwasawa \(\mu\)-invariant is \(0\) on the single-oscillator
component, \Cref{v4-arith:pos:cor:mu-zero-heisenberg}).
\end{definition}

\begin{remark}[the archimedean operator requirement]
\label{v4-arith:prim-zeta:rem:consistency-vb}
Finite occupation energy is defined on the algebraic primary tower,
and its powers have the usual maximal weighted sequence domains in its
Hilbert completion. A Gaussian test line does not carry a continuous
energy spectrum. An archimedean dilation operator requires an enlarged
function space and a specified domain; the Gaussian is not itself a
dilation eigenvector. Consequently a global descendant or Virasoro
operator cannot be defined from the mixed character alone.
\end{remark}

\begin{remark}[determinant rank and Virasoro reconciliation with
\Cref{v4-arith:rh:conj:primary-Virasoro}]
\label{v4-arith:prim-zeta:rem:Virasoro-reconciliation}
On the primary subspace
\(\mathcal{V}^{\mathrm{prim}}_{\mathrm{Heis},\mathrm{prim}}\), the
finite conformal modes \(m\geq -1\) act on finite-support vectors. The
negative modes \(m\leq -2\) do not form global sums on the restricted
vacuum, so a Virasoro action is not produced by the character theorem.
The value proved in \Cref{v4-arith:thm:regularised-c} is the
determinant rank residue \(r_{\zeta}=2\), obtained from the pole of
\(\zeta(s)\) at \(s=1\). The conjectural primary-Virasoro bracket
requires a separate operator construction whose central charge matches
this residue.
\end{remark}

\section{Character on Primary via Mellin--Theta}
\label{v4-arith:prim-zeta:character-primary}

\subsection{Finite-Place Primary Character}
\label{v4-arith:prim-zeta:finite-char}

\begin{theorem}[finite-place primary character]
\label{v4-arith:prim-zeta:thm:finite-char}
\ClaimStatusProvedHere. The character of
\(\bigotimes'_{p}\mathcal{H}^{(p)}_{\mathrm{prim}}\) under the
HY-zeta grading is the Dirichlet series \(\zeta(s)\), convergent
absolutely for \(\mathrm{Re}(s)>1\) and admitting meromorphic
continuation to \(\mathbb{C}\) with a simple pole at \(s=1\) and no
other poles:
\begin{equation}
\label{v4-arith:prim-zeta:eq:finite-char}
\mathrm{char}\Bigl(\bigotimes_{p\text{ prime}}^{\prime}\mathcal{H}^{(p)}_{\mathrm{prim}}\Bigr)(s)
\;=\;\prod_{p\text{ prime}}\bigl(1-p^{-s}\bigr)^{-1}
\;=\;\zeta(s),
\qquad\mathrm{Re}(s)>1.
\end{equation}
\end{theorem}

\begin{proof}
Per prime, the complex trace $\operatorname{Tr}(e^{-sH_p})$
is the absolutely convergent geometric series
$\sum_{n\geq0}p^{-ns}$ when $\operatorname{Re}s>0$.
A restricted tensor basis is indexed by positive integers through
unique factorization. Its energy is $\log n$, hence its trace is
$\sum_{n\geq1}n^{-s}$. For $\operatorname{Re}s>1$, absolute
convergence permits the regrouping into prime powers and yields
the Euler product. This argument concerns an unnormalized operator
trace. A KMS state requires a real inverse temperature and the
normalization in \Cref{v4-arith:prim-zeta:prop:BC-ergodic}.

Convergence: the product \(\prod_p(1-p^{-s})^{-1}\) converges absolutely for
\(\mathrm{Re}(s)>1\) by the bound \(|\log(1-p^{-s})^{-1}|\leq\sum_{k\geq 1}p^{-k\sigma}/k\leq 2p^{-\sigma}\)
for \(\mathrm{Re}(s)\geq 1\), and \(\sum_p p^{-\sigma}\) converges for
\(\sigma=\mathrm{Re}(s)>1\) (this is the prime zeta function \(P(\sigma)\),
finite for \(\sigma>1\); Landau 1909). Meromorphic continuation: the identity
\(\zeta(s)=(1-2^{1-s})^{-1}\sum_{n\geq 1}(-1)^{n-1}n^{-s}\)
(alternating zeta) gives continuation to \(\mathrm{Re}(s)>0\), then
functional equation gives continuation to \(\mathbb{C}\). The simple
pole at \(s=1\) is a standard computation
(\Cref{v4-arith:thm:regularised-c}: \(\mathrm{Res}_{s=1}\zeta(s)=1\)).

Julia's 1990 primon gas identification
(\Cref{v4-arith:thm:P1-resolution} Step 1): the per-prime
single-oscillator Fock partition function equals
\((1-p^{-\beta})^{-1}\) at inverse temperature \(\beta\), and the
primon gas total partition is \(\zeta(\beta)\), identifying the
HY-zeta character with the Dirichlet series identically.
\end{proof}

\subsection{Archimedean Primary Character}
\label{v4-arith:prim-zeta:arch-char}

\begin{theorem}[archimedean primary character]
\label{v4-arith:prim-zeta:thm:arch-char}
\ClaimStatusProvedHere. For \(\operatorname{Re}s>0\), the Mellin character of the specified
Gaussian spanning \(\mathcal{H}^{(\infty)}_{\mathrm{prim}}\) is
\begin{equation}
\label{v4-arith:prim-zeta:eq:full-arch-char}
\mathrm{char}\bigl(\mathcal{H}^{(\infty)}_{\mathrm{prim}}\bigr)(s)
\;=\;\int_{\mathbb{R}}e^{-\pi x^{2}}|x|^{s-1}dx
\;=\;\pi^{-s/2}\Gamma(s/2)
\;=\;\Gamma_{\mathbb{R}}(s),
\end{equation}
the archimedean local \(L\)-factor of the trivial Dirichlet character.
\end{theorem}

\begin{proof}
\Cref{v4-arith:prim-zeta:prop:arch-char}: the Mellin pairing of the
Gaussian primary state with \(|x|^{s-1}\) gives \(\Gamma_{\mathbb{R}}(s)\).
The Mellin--theta machinery of Tate 1950 \S4.4 identifies this with
the archimedean local \(L\)-factor: \(L_{\infty}(s,\chi_{\mathrm{triv}})=\Gamma_{\mathbb{R}}(s)\),
with the Gaussian and its amplitude fixed as in the definition.
\end{proof}

\subsection{Adelic Character: \(\Lambda_{\zeta}(s)\)}
\label{v4-arith:prim-zeta:adelic-char}

\begin{theorem}[primitive zeta-character identity]
\label{v4-arith:prim-zeta:thm:adelic-char}
\ClaimStatusProvedHere. The character of the primary subspace
\(\mathcal{V}^{\mathrm{prim}}_{\mathrm{Heis},\mathrm{prim}}\) with the mixed operation of \Cref{v4-arith:prim-zeta:mixed-character}
is the meromorphic Tate character:
\begin{equation}
\label{v4-arith:prim-zeta:eq:adelic-char}
\mathrm{char}\bigl(\mathcal{V}^{\mathrm{prim}}_{\mathrm{Heis},\mathrm{prim}}\bigr)(s)
\;=\;\Gamma_{\mathbb{R}}(s)\cdot\zeta(s)\;=\;\Lambda_{\zeta}(s).
\end{equation}
\end{theorem}

\begin{proof}
Factorise the adelic primary as
\(\mathcal{V}^{\mathrm{prim}}_{\mathrm{Heis},\mathrm{prim}}=\mathcal{H}^{(\infty)}_{\mathrm{prim}}\otimes_{\mathbb{C}}\bigotimes'_{p}\mathcal{H}^{(p)}_{\mathrm{prim}}\)
(\Cref{v4-arith:prim-zeta:def:primary-subspace}). The character of
a tensor product is the product of characters:
\begin{equation}
\mathrm{char}(\mathcal{V}^{\mathrm{prim}}_{\mathrm{Heis},\mathrm{prim}})(s)
\;=\;\mathrm{char}(\mathcal{H}^{(\infty)}_{\mathrm{prim}})(s)\cdot\mathrm{char}\bigl(\bigotimes'_{p}\mathcal{H}^{(p)}_{\mathrm{prim}}\bigr)(s).
\end{equation}
By \Cref{v4-arith:prim-zeta:thm:arch-char},
\(\mathrm{char}(\mathcal{H}^{(\infty)}_{\mathrm{prim}})(s)=\Gamma_{\mathbb{R}}(s)\).
By \Cref{v4-arith:prim-zeta:thm:finite-char},
\(\mathrm{char}(\bigotimes'_{p}\mathcal{H}^{(p)}_{\mathrm{prim}})(s)=\zeta(s)\).
Product: \(\Gamma_{\mathbb{R}}(s)\zeta(s)=\Lambda_{\zeta}(s)\), the standard
completion by Tate 1950 \S4.4 / Cassels--Fr\"ohlich 1967 \S13.4.
\end{proof}

\begin{corollary}[primary character of Heisenberg is \(\Lambda_{\zeta}\)]
\label{v4-arith:prim-zeta:cor:primary-is-xi}
\ClaimStatusProvedHere.
\(\mathrm{char}(\mathcal{V}^{\mathrm{prim}}_{\mathrm{Heis},\mathrm{prim}})(s)=\Lambda_{\zeta}(s)\),
a meromorphic function of \(s\) with simple poles only at \(s=0\) and
\(s=1\) (both residues \(\pm 1\)), and with functional equation
\(\Lambda_{\zeta}(s)=\Lambda_{\zeta}(1-s)\) (Riemann 1859; Tate 1950).
\end{corollary}

\begin{proof}
Combine \Cref{v4-arith:prim-zeta:thm:adelic-char} with Tate's
functional equation
(\Cref{v4-arith:thm:P2-resolution}(3)).
\end{proof}

\section{Koszul-Dual Character (Bost--Connes Side)}
\label{v4-arith:prim-zeta:koszul-dual-char}

\subsection{Cobar of Primary Heisenberg}
\label{v4-arith:prim-zeta:cobar-primary}

\begin{definition}[native bar and dual carriers]
\label{v4-arith:prim-zeta:def:cobar-primary}
Write \(C_{\mathrm{prim}}\) for the native arithmetic bar object when
its restriction to the primary generators is constructed. Write
\(A_{\mathrm{prim}}=C_{\mathrm{prim}}^\vee\) only after its coefficient
field, grading, topology and continuous-dual convention have been fixed.
An identification of this dual with a cobar object is a comparison map,
not a consequence of the notation. Identifying its trace-carrying sector
with the occupation basis and reading that trace at \(1-s\) is part of
\Cref{v4-arith:prim-zeta:native-character-hypothesis}. The stronger
Bost--Connes algebra presentation requires compatible multiplication
and phase maps in addition.
\end{definition}

\subsection{Bost--Connes KMS Partition on Primary}
\label{v4-arith:prim-zeta:BC-KMS}

\begin{theorem}[operator partition and finite primary character]
\label{v4-arith:prim-zeta:thm:BC-partition}
\ClaimStatusProvedHere. On the occupation Hilbert space and with
$Hv_n=\log n\,v_n$, the holomorphic trace is
\begin{equation}
 \label{v4-arith:prim-zeta:eq:BC-partition}
 Z(s)=\operatorname{Tr}(e^{-sH})
     =\prod_p(1-p^{-s})^{-1}=\zeta(s),
       \qquad\operatorname{Re}s>1.
\end{equation}
For real $s=\beta>1$, this is the Gibbs partition trace of the
represented arithmetic shift system. Its normalized state is
\Cref{v4-arith:prim-zeta:global-shift-state}. Meromorphic
continuation of the scalar $\zeta(s)$ does not continue this
formula as a nuclear operator trace outside its convergence domain.
\end{theorem}

\begin{proof}
The occupation basis gives the diagonal entries $n^{-s}$.
Their absolute values sum to $\zeta(\operatorname{Re}s)$, finite
precisely when $\operatorname{Re}s>1$. Unique factorization
regroups this absolutely convergent sum into its prime geometric
series, exactly as in \Cref{v4-arith:prim-zeta:thm:finite-char}.
For real $\beta$, division by $Z(\beta)$ gives a positive
normalized functional. These constructions identify the displayed
operator partition and the finite primary character; they do not
identify the arithmetic cobar algebra with an operator algebra.
Nor do they produce a cyclic functional on ordinary Hochschild
homology. That comparison uses the specified twisted or nuclear
coefficients of \Cref{v4-arith:thermal:thm:density-chain-map}.
\end{proof}

\subsection{Archimedean Factor via Fourier Duality}
\label{v4-arith:prim-zeta:arch-fourier}

\begin{proposition}[conditional dual archimedean character]
\label{v4-arith:prim-zeta:prop:arch-cobar}
Under \Cref{v4-arith:prim-zeta:native-character-hypothesis}, the native
dual archimedean functional has value
\begin{equation}
 \Gamma_{\mathbb R}(1-s)=\pi^{-(1-s)/2}\Gamma((1-s)/2).
 \label{v4-arith:prim-zeta:eq:arch-cobar}
\end{equation}
\end{proposition}
\begin{proof}
The assumed map identifies that functional with the Gaussian Mellin
integral at \(1-s\). The same change of variable as before evaluates
it when \(\operatorname{Re}s<1\), then its scalar meromorphic
continuation extends the formula. The Gaussian's Fourier self-duality
alone does not construct the native map.
\end{proof}

\subsection{Adelic dual character}
\label{v4-arith:prim-zeta:koszul-dual-adelic}
\begin{theorem}[conditional native dual trace]
\label{v4-arith:prim-zeta:thm:koszul-dual-adelic}
Under \Cref{v4-arith:prim-zeta:native-character-hypothesis}, with the
stated compatible assembly of its trace-carrying sectors,
\begin{equation}
 \mathrm{char}_{A^!}(s)=\Gamma_{\mathbb R}(1-s)\zeta(1-s)
                     =\Lambda_\zeta(1-s).
 \label{v4-arith:prim-zeta:eq:koszul-dual-adelic}
\end{equation}
\end{theorem}
\begin{proof}
The two assumed trace comparisons identify the finite and archimedean
functionals with the explicitly computed factors. They converge
simultaneously when \(\operatorname{Re}s<0\); their product has the
stated meromorphic continuation. No KMS state with complex inverse
temperature is used, and no native multiplication is inferred.
\end{proof}

\section{Primitive Zeta-Character Theorem}
\label{v4-arith:prim-zeta:main-theorem}

\subsection{Assembly}
\label{v4-arith:prim-zeta:assembly}

\begin{proof}[Proof of \Cref{v4-arith:prim-zeta:thm:main}]
The character identity
(\ref{v4-arith:prim-zeta:eq:main})
\begin{equation*}
\mathrm{char}(\mathcal{V}^{\mathrm{prim}}_{\mathrm{Heis},\mathrm{prim}})(s)
\;=\;\Gamma_{\mathbb{R}}(s)\,\zeta(s)\;=\;\Lambda_{\zeta}(s)
\end{equation*}
is \Cref{v4-arith:prim-zeta:thm:adelic-char}: factorise the adelic
primary into finite-place and archimedean components, apply
\Cref{v4-arith:prim-zeta:thm:finite-char} (Julia 1990 primon gas
+ Euler product) and
\Cref{v4-arith:prim-zeta:thm:arch-char} (Tate 1950 Mellin--theta on
Gaussian), tensor to obtain
\(\Gamma_{\mathbb{R}}(s)\zeta(s)=\Lambda_{\zeta}(s)\).

The character identity
(\ref{v4-arith:prim-zeta:eq:koszul-dual})
is \Cref{v4-arith:prim-zeta:thm:koszul-dual-adelic}: the Koszul-dual
character satisfies \(\mathrm{char}_{A^{!}}(s)=\Lambda_{\zeta}(1-s)\), so
\(\mathrm{char}_{C}(s)=\mathrm{char}_{A^{!}}(1-s)=\Lambda_{\zeta}(1-s)\big|_{1-s\mapsto s}=\Lambda_{\zeta}(s)\).
Tate's functional equation \(\Lambda_{\zeta}(s)=\Lambda_{\zeta}(1-s)\)
(\Cref{v4-arith:thm:P2-resolution}(3)) then gives
\(\mathrm{char}_{C}(s)=\Lambda_{\zeta}(s)=\Lambda_{\zeta}(1-s)=\mathrm{char}_{A^{!}}(s)\)
as an identity of meromorphic functions.

The mixed scalar identity is unconditional. The native character
step uses \Cref{v4-arith:prim-zeta:native-character-hypothesis} and its
trace and assembly maps. The arithmetic pairing and operations require
further compatible constructions.
\end{proof}

\section{Native character and pairing requirements}
\label{v4-arith:prim-zeta:consequence-p1}
\label{v4-arith:prim-zeta:reduction-sub-items}

The finite thermal representation proves its own trace and state
formulas. The Gaussian calculation proves its Mellin integral. Their
product and the functional equation are identities of scalar functions.
The original arithmetic character target additionally requires the
native bar and cobar trace sectors and comparison maps of
\Cref{v4-arith:prim-zeta:native-character-hypothesis}.

\begin{corollary}[conditional native character transfer]
\label{v4-arith:prim-zeta:cor:p1-residual-reduced}
Under the native character realization hypothesis, the mixed scalar
calculation gives the character part of the arithmetic duality target.
It gives no construction of continuous duals, derived completion,
faithful lifting or a multiplicative native comparison.
\end{corollary}
\begin{proof}
The assumed trace maps identify the computed scalar with the native
character by definition of trace preservation. Their existence and all
stronger categorical properties remain separate hypotheses.
\end{proof}

\begin{proposition}[conditional functional equation for native characters]
\label{v4-arith:prim-zeta:fingerprint}
\label{v4-arith:prim-zeta:prop:char-functional-eq}
Under those trace maps,
\begin{equation}
 \mathrm{char}_{C}(s)=\mathrm{char}_{A^!}(s)
 \quad\Longleftrightarrow\quad
 \Lambda_\zeta(s)=\Lambda_\zeta(1-s).
 \label{v4-arith:prim-zeta:eq:char-functional-eq}
\end{equation}
\end{proposition}
\begin{proof}
Substitute the two character values identified by the hypothesis.
Without those identifications, the scalar functional equation is not a
native Koszul character theorem.
\end{proof}

\begin{remark}[trace and algebra comparisons]
\label{v4-arith:prim-zeta:disjoint-sources}
\label{v4-arith:prim-zeta:rem:sources}
\label{v4-arith:prim-zeta:rem:sources-disjoint}
\label{v4-arith:prim-zeta:rem:p1-form}
\label{v4-arith:prim-zeta:rem:caveat-fingerprint}
Even equality of characters does not identify multiplication, higher
operations or the bar--cobar pairing. Continuous duals, derived
completion and faithful lifting must be constructed on their stated
coefficient and topological carriers. The character part of a native
duality claim retains its trace-map hypothesis.
\end{remark}

\subsection{The first native realization map}
\label{v4-arith:prim-zeta:remaining-four}
\label{v4-arith:prim-zeta:residual-table}
\label{v4-arith:prim-zeta:tab:residual}
\label{v4-arith:prim-zeta:summary}
The phase obstruction and twisted Hochschild calculation of
\Cref{v4-arith:thermal:chapter} give explicit tests for the first native
realization map. The local normalized state does not descend to ordinary
Hochschild degree zero, whereas the density multiplication map gives a
chain map with its stated twisted coefficient module. Whether the
native arithmetic trace uses that coefficient module, and how it maps
to a Hall or BPS construction, remains a separate realization problem.
